% Version prepared for author review on 21 August 2026.
% Springer Nature LaTeX authoring template v3.1 (December 2024).
% Target: npj Unconventional Computing, Article.
% Lines containing "AUTHOR CHECK" require verification before submission.

\documentclass[referee,lineno,pdflatex,sn-nature]{sn-jnl}

\usepackage{graphicx}
\usepackage{amsmath,amssymb,amsfonts}
\usepackage{bm}
\usepackage{booktabs}
\usepackage{array}
\usepackage{mathrsfs}
\usepackage{xcolor}
\usepackage{textcomp}

\raggedbottom

\newcommand{\kBT}{k_{\mathrm B}T}
\newcommand{\vect}[1]{\boldsymbol{#1}}
\newcommand{\Ehat}{\widehat E}
\newcommand{\Fhat}{\widehat{\vect F}}
\newcommand{\authorcheck}[1]{\textbf{[AUTHOR CHECK: #1]}}

\begin{document}

\title[Local learning for physical generative models]
{Local learning rules for nonequilibrium physical generative models}

\author*[1]{\fnm{Cyrill} \sur{B\"osch}}\email{cb7454@princeton.edu}
\author[1]{\fnm{Geoffrey} \sur{Roeder}}
\author[2]{\fnm{Marc} \sur{Serra-Garcia}}
\author[1]{\fnm{Ryan P.} \sur{Adams}}

\affil*[1]{\orgdiv{Department of Computer Science},
  \orgname{Princeton University},
  \orgaddress{\city{Princeton}, \postcode{08540}, \state{NJ},
  \country{USA}}}
\affil[2]{\orgname{AMOLF},
  \orgaddress{\street{Science Park 104}, \postcode{1098 XG},
  \city{Amsterdam}, \country{The Netherlands}}}

\abstract{Score-based generative models produce samples through a time-dependent
nonequilibrium stochastic process, but their neural implementations are usually
trained by global backpropagation. We derive a parameter-local rule for training
the required score field in a conservative physical system. We parameterize a
time-dependent score energy in a network of nonlinear overdamped oscillators. For this
model, we derive an estimator whose expectation equals the score-matching gradient.
It uses local force components measured with the full state clamped and analytic
derivatives of the corresponding energy feature. When forces cannot be measured, a
one-step contrastive estimator approaches the same gradient as the release time
tends to zero. We illustrate the construction
numerically for a two-mode distribution and, as a qualitative proof of principle,
for grayscale MNIST images of zeros and ones using a $28\times28$ oscillator network.
This connects local physical learning with nonequilibrium generative sampling.}

\keywords{physical learning, score matching, diffusion models, thermodynamic
computing, nonlinear oscillators, probabilistic computing}

\maketitle

\section{Introduction}\label{sec:introduction}

Physical computing maps a computation onto a device's native dynamics
\cite{jaeger2023toward,finocchio2024roadmap,
melanson2025thermodynamic,aifer2024thermodynamic,duffield2025thermodynamic,
wanjura2024fully,lin2025memristive,nakajima2015information}. The physical parameters
must still be trained for the target task. Physical learning derives updates from
measurements local to each tunable element
\cite{momeni2024training,stern2023learning,stern2021supervised,
dillavou2022demonstration,dillavou2024machine,altman2024experimental}. Many such
rules use equilibria or steady states of the physical dynamics
\cite{scellier2017equilibrium,baldi1991contrastive,laydevant2024training,
guzman2025unsupervised}.

Diffusion models instead use explicitly time-dependent dynamics. A score-based
generative model (SGM) evolves a simple distribution into a target distribution
through a finite stochastic process \cite{song2019generative,
song2021scorebased}. The reverse dynamics depend on the score, the spatial gradient
of the log density at every noise time. Neural score models are normally trained
by differentiating a digital model. A physical realization should instead learn the
required force field without differentiating a numerical model of the device.

Recent work has mapped generative diffusion inference to physical systems. Passive
diffractive optical systems have been programmed to perform denoising inference
\cite{oguz2024optical}, and analog fixed-point formulations retain conventional
model training while mapping inference to noisy optical hardware
\cite{chu2026analog}. A resistive-memory neural differential-equation solver has
implemented time-continuous score-based diffusion with an offline-trained analog
neural network and feedback integrator \cite{yang2026resistive}. Denoising
thermodynamic models are built from sparse, locally
connected Boltzmann machines and have been analysed together with an all-transistor
probabilistic architecture \cite{jelincic2026efficient}. These studies either
perform or model physical diffusion inference and, in some cases, provide device-level
performance projections. We focus on training a time-indexed physical score energy
from local observables.

Local training has also been developed through equilibrium propagation,
contrastive procedures, and related steady-state constructions
\cite{scellier2017equilibrium,stern2021supervised,dillavou2022demonstration}.
Related work uses equilibrium propagation to obtain local denoising-score
gradients from a zero-nudge limit and multiple equilibrium states in a bilinear
thermodynamic model \cite{de2026symmetric}. Our direct rule instead uses one fully
clamped force configuration for each sample at a positive noise time. The
expectation of the resulting local estimator equals the population score-matching
gradient. A short-time contrastive release provides an alternative when the force
is not directly observable.
Here, locality means that each parameter update uses only the coordinates and force
signals associated with its energy feature.

We make three contributions. First, we express the score through a time-dependent
learned energy and derive a local force-matching rule for its coefficients.
All per-sample coefficient-gradient contributions can be accumulated in parallel
from the same clamped configuration, provided that the corresponding force
components are readable.
Second, we formulate a one-step contrastive semi-gradient that converges to the same
score-matching gradient as the release duration tends to zero. Third, we map the
learned score to the open-loop energy protocol of a nonlinear overdamped oscillator
network and demonstrate the construction numerically on a bimodal distribution and
on grayscale MNIST images. The numerical examples illustrate the learning
construction; they do not demonstrate a fabricated device or compare end-to-end
energy or speed.

\section{Results}\label{sec:results}

\subsection{Reverse diffusion as driven Langevin dynamics}

Generative modelling seeks samples from a target distribution $p$ observed only
through a dataset
$\mathcal D=\{\vect x^{(m)}\in\mathbb R^N\}_{m=1}^{M}$. We use
$t\in[0,\tau]$ for forward time. The forward Ornstein--Uhlenbeck process is
\begin{equation}\label{eq:forward}
  \mathrm d\vect X_t=-\vect X_t\,\mathrm dt
  +\sqrt{2\kBT}\,\mathrm d\vect W_t,
\end{equation}
with $p_0=p$. Its transition law is
\begin{equation}\label{eq:ou-transition}
  q_t(\vect x_t\mid\vect x_0)
  =\mathcal N(a_t\vect x_0,\sigma_t^2\vect I),\qquad
  a_t=e^{-t},\quad \sigma_t^2=\kBT(1-e^{-2t}).
\end{equation}
For the experiments, the dataset defines the forward distribution
\begin{equation}\label{eq:empirical-smoothed}
 p_t^{\mathcal D}(\vect x)=\frac{1}{M}\sum_{m=1}^{M}
 q_t(\vect x\mid\vect x^{(m)}), \qquad t>0.
\end{equation}
For sufficiently large $\tau$, $p_\tau$ is approximated by
$\mathcal N(\vect 0,\kBT\vect I)$.

Let $s\in[0,\tau]$ denote reverse time, with $s=0$ at the noisy endpoint. The
reverse process is \cite{anderson1982reverse,song2021scorebased}
\begin{equation}\label{eq:reverse}
  \mathrm d\vect Y_s=
  \left[\vect Y_s+2\kBT\nabla_{\vect x}\log p_{\tau-s}(\vect Y_s)\right]
  \mathrm ds+\sqrt{2\kBT}\,\mathrm d\overline{\vect W}_s,
\end{equation}
where $\vect Y_0\sim p_\tau$. With an exact score and this initial distribution,
$\vect Y_\tau\sim p_0$.

We parameterize the score at forward time $t$ through a learned score energy,
\begin{equation}\label{eq:force-score}
  \Fhat_{\vect\theta,t}(\vect x)
  =-\nabla_{\vect x}\Ehat_{\vect\theta,t}(\vect x),\qquad
  \vect s_{\vect\theta,t}(\vect x)
  =\frac{\Fhat_{\vect\theta,t}(\vect x)}{\kBT}.
\end{equation}
The energy replayed by the physical system at sampling time $s$ is
\begin{equation}\label{eq:reverse-energy}
  U_{\vect\theta,s}(\vect x)
  =2\Ehat_{\vect\theta,\tau-s}(\vect x)
  -\frac{1}{2}\lVert\vect x\rVert^2 .
\end{equation}
Its overdamped force is
$-\nabla U_{\vect\theta,s}=\vect x+2\Fhat_{\vect\theta,\tau-s}$,
which reproduces the learned drift in Eq.~\eqref{eq:reverse}. The score energy and
the physical reverse energy are denoted by $\Ehat$ and $U$, respectively.

\subsection{Polynomial oscillator model}

We represent $\Ehat_{\vect\theta,t}$ using a low-order polynomial energy on an
interaction graph $\mathcal E$. We assign a fixed ordering to each retained pair
$(n,m)$ because the $\chi_{nm}$ and $\omega_{nm}$ features are
exchange-asymmetric. The on-site contribution is
\begin{equation}\label{eq:local-energy}
 \Ehat^{\mathrm l}_{\vect\theta,t}(\vect x)
 =\sum_{n=1}^{N}\left[
 \frac{\alpha_n(t)}{2}x_n^2+
 \frac{\beta_n(t)}{4}x_n^4+
 \frac{\gamma_n(t)}{6}x_n^6+f_n^{\mathrm{ext}}(t)x_n\right],
\end{equation}
and the pair contribution for $(n,m)\in\mathcal E$ is
\begin{align}\label{eq:pair-energy}
 \Ehat^{\mathrm{pair}}_{nm,t}(x_n,x_m)
 ={}&\frac{\kappa_{nm}(t)}{2}(x_n-x_m)^2
 +\frac{\lambda_{nm}(t)}{4}(x_n-x_m)^4 \nonumber\\
 &+\chi_{nm}(t)x_nx_m^2+\omega_{nm}(t)x_n^2x_m .
\end{align}
The complete learned energy is
$\Ehat_{\vect\theta,t}=\Ehat^{\mathrm l}_{\vect\theta,t}
+\sum_{(n,m)\in\mathcal E}\Ehat^{\mathrm{pair}}_{nm,t}$.
We constrain $\gamma_n(t)>0$ to keep the learned energy coercive.

Oscillator coordinates provide continuous state variables. On-site nonlinearities
and pair interactions can represent the polynomial energy, while thermal
fluctuations and clamped force readout supply the noise and training signal. The
$\chi_{nm}$ and $\omega_{nm}$ features are exchange-asymmetric polynomial monomials.
Because they belong to a scalar energy, the resulting force field remains
conservative. Related monomials can arise in optomechanical, magnetic, and
geometrically nonlinear systems
\cite{sankey2010strong,serra2018tunable,serra2016mechanical}.

\subsection{Force matching gives parameter-local score gradients}

At a fixed positive noise time, let the learned energy be linear in features,
\begin{equation}\label{eq:feature-energy}
  \Ehat_{\vect\theta}(\vect x)=\sum_i\theta_i\phi_i(\vect x),
\end{equation}
and let $\mathcal S_i$ contain the coordinates on which $\phi_i$ depends. The score-matching objective is \cite{hyvarinen2005scorematching}
\begin{equation}\label{eq:sm-objective}
 J_{\mathrm{SM}}(\vect\theta;t)=
 \mathbb E_{\vect X\sim p_t}\left[
 \nabla_{\vect x}\!\cdot\vect s_{\vect\theta}(\vect X)
 +\frac{1}{2}\lVert\vect s_{\vect\theta}(\vect X)\rVert^2\right].
\end{equation}
Using Eq.~\eqref{eq:force-score}, its derivative is
\begin{equation}\label{eq:local-gradient}
 \frac{\partial J_{\mathrm{SM}}}{\partial\theta_i}
 =\mathbb E_{p_t}\left[
 -\frac{\Delta\phi_i(\vect X)}{\kBT}
 -\frac{\nabla\phi_i(\vect X)\cdot\Fhat_{\vect\theta}(\vect X)}
 {(\kBT)^2}\right].
\end{equation}
Equation~\eqref{eq:local-gradient} is exact for the
objective~\eqref{eq:sm-objective}. In the experiments, we replace $p_t$ by
$p_t^{\mathcal D}$ from Eq.~\eqref{eq:empirical-smoothed}; the same exactness
statement holds at $t>0$. Finite samples and imperfect force measurements
introduce estimation error. Each sample requires preparing and clamping the full
network state. Because $\nabla\phi_i$ and $\Delta\phi_i$ vanish outside $\mathcal S_i$,
the contribution for $\theta_i$ uses only the coordinates and force components in
$\mathcal S_i$. The curvature term is an analytic function of those coordinates;
the other term is obtained from the force of the clamped system.

For the on-site Duffing feature $\phi_{\beta_n}=x_n^4/4$,
\begin{equation}\label{eq:duffing-gradient}
 \frac{\partial J_{\mathrm{SM}}}{\partial\beta_n}
 =\mathbb E_{p_t}\left[-\frac{3X_n^2}{\kBT}
 -\frac{X_n^3\widehat F_{\vect\theta,n}(\vect X)}{(\kBT)^2}\right].
\end{equation}
For a coupling feature $\phi_i(x_n,x_m)$, the update uses $x_n$, $x_m$, and the
corresponding components $\widehat F_n$ and $\widehat F_m$ of the total learned
force. All per-sample coefficient-gradient contributions can therefore be
accumulated in parallel from one clamped configuration.

The learned force can be measured directly if the training configuration implements
$\Ehat$. If instead the complete reverse energy $U$ is active, clamping the network
and measuring $\vect F^{\mathrm{rev}}=-\nabla U$ gives
$\Fhat=(\vect F^{\mathrm{rev}}-\vect x)/2$. No differentiation through a numerical
device model is required.

\subsection{Short-time contrastive estimator}

When force components cannot be measured, a short physical release provides an
alternative. Starting from $\vect X_t\sim p_t$, we approximate a Langevin release
of duration $\delta$ under the learned score energy by one Euler--Maruyama step:
\begin{equation}\label{eq:cd-release}
 \vect X_t^- =\vect X_t-\delta\nabla_{\vect x}\Ehat_{\vect\theta}(\vect X_t)
 +\sqrt{2\delta\kBT}\,\vect W,
 \qquad \vect W\sim\mathcal N(\vect0,\vect I).
\end{equation}
For a linearly parameterized energy, define the stop-gradient semi-gradient
\begin{equation}\label{eq:cd-semigradient}
 g_i^{\mathrm{CD1}}(\vect\theta;t,\delta)
 =\mathbb E\left[\phi_i(\vect X_t^-)-\phi_i(\vect X_t)\right].
\end{equation}
In hardware, the corresponding state can be obtained by running the continuous
Langevin dynamics under $\Ehat_{\vect\theta}$ for duration $\delta$; the displayed
step approximates those dynamics to first order in $\delta$. The full coupled state
evolves during the release. Locality applies to the feature difference used for a
parameter update, not to the release dynamics. The update is formed directly from
this measured feature difference; no derivative of the release trajectory is
taken. Under the smoothness assumptions of the
score-matching--contrastive-divergence expansion \cite{hyvarinen2007connections},
\begin{equation}\label{eq:cd-limit}
 -\frac{g_i^{\mathrm{CD1}}}{\delta(\kBT)^2}
 =\frac{\partial J_{\mathrm{SM}}}{\partial\theta_i}+O(\delta).
\end{equation}
Repeated releases reduce Monte Carlo variance but not the finite-$\delta$ bias. For
$\phi_{\beta_n}=x_n^4/4$, Eq.~\eqref{eq:cd-semigradient} is the average of
$[(X_{t,n}^-)^4-X_{t,n}^4]/4$. This operation needs the feature definition,
temperature, and measured states.


\subsection{Nonequilibrium sampling of a bimodal distribution}

We first consider a two-dimensional mixture of two Gaussians,
\begin{equation}
 p(x_1,x_2)=\sum_{j=1}^{2}\varpi_j
 \mathcal N\!\left((x_1,x_2)\mid\vect\mu_j,\vect\Sigma_j\right).
\end{equation}
\authorcheck{Insert the mixture weights, mean vectors, covariance matrices, and
number of training samples.}
We train a two-oscillator network using Eq.~\eqref{eq:local-gradient} and generate
$S=2000$ independent trajectories of duration $\tau$. Figure~\ref{fig:mog}
shows that the samples populate both modes and that their relative frequencies
visually agree with the target weights on this time scale.

For an equilibrium comparison with the same endpoint energy, we use
$\Ehat_{\vect\theta,0}$ as a static Boltzmann energy. One persistent Langevin chain,
without restarts, is evolved for total time $S\tau$ and recorded every $\tau$. On the simulated time
scale it remains in the right basin even though the learned energy contains both
basins. For this initialization and runtime, the persistent chain mixes slowly
because transitions between the basins are rare. Each reverse trajectory instead
reaches a mode through a finite-time denoising path.

\begin{figure}[t]
\centering
\includegraphics[width=\textwidth]{fig1.png}
\caption{\textbf{Two-dimensional Gaussian-mixture example.} Top row: learned
nonequilibrium reverse process. Bottom row: the matched persistent equilibrium
chain. (a, d) Histograms of $2000$ recorded samples and exact marginal densities.
(b, e) Learned endpoint score energy $\Ehat_{\vect\theta,0}$ with samples.
(c, f) Boltzmann density defined by the same endpoint score energy; yellow dashed
curves show target-density contours. The nearly uniform dark-blue background reflects
the linear colour scale and broad dynamic range of the density.}
\label{fig:mog}
\end{figure}

The CD1 estimator learns a different energy and visually captures both modes in
this example (Fig.~\ref{fig:cd1}). The force-matching
estimator is used for the larger example below as it avoids the $O(\delta)$ bias.

\begin{figure}[t]
\centering
\includegraphics[width=\textwidth]{fig3.png}
\caption{\textbf{Two-dimensional example trained with the CD1 estimator.}
Top row: learned reverse-process samples, learned score energy, and the associated
endpoint Boltzmann density. Bottom row: matched persistent equilibrium comparison.
Histograms contain $2000$ recorded samples; yellow dashed curves show target-density
contours. The update uses the feature differences defined in
Eq.~\eqref{eq:cd-semigradient}.}
\label{fig:cd1}
\end{figure}

\subsection{MNIST generation with 784 oscillators}

We next train the polynomial oscillator model on grayscale MNIST images restricted
to digits 0 and 1. Each of the $28\times28=784$ pixels is represented by one
continuous oscillator coordinate. The graph used in the reported experiment
contains couplings between positions separated by as many as 14 sites.
Figure~\ref{fig:mnist} compares representative data images with samples from the
learned reverse protocol. The generated patterns are recognizable as zeros or ones,
but are blurrier and less regular than the data examples.

\begin{figure}[t]
\centering
\includegraphics[width=\textwidth]{fig2.png}
\caption{\textbf{MNIST samples.} (a) Representative grayscale images of zeros and
ones from the training data. (b) Endpoint samples generated by the locally trained
$28\times28$ oscillator network.}
\label{fig:mnist}
\end{figure}

\subsection{System boundary and scaling}

The proposed computer is hybrid. The physical system generates the stochastic
dynamics and nonlinear forces. An external controller initializes the state and
replays the learned schedule $\vect\theta^*(\tau-s)$, using the same schedule for
every sample. The controller and its programming cost are therefore part of the
system. Training also requires preparing and clamping corrupted data states,
measuring forces or running short releases, and updating the coefficients. The CD1
releases use the learned score energy $\Ehat$ rather than the reverse energy $U$.

For the basis above, the number of programmed coefficients is
\begin{equation}\label{eq:pcount}
 P_{\mathrm{elem}}=4N+4|\mathcal E|.
\end{equation}
A schedule stored at $N_t$ noise times contains $P_{\mathrm{elem}}N_t$ values.

Let $K_{\mathrm{opt}}$ be the number of optimizer updates per noise time, $B$ the
batch size, and $R$ the number of CD1 releases per clamped state. Force matching uses
$O(N_tK_{\mathrm{opt}}B)$ clamped configurations and
$O(P_{\mathrm{elem}}N_tK_{\mathrm{opt}}B)$ gradient contributions. CD1 uses
$O(RN_tK_{\mathrm{opt}}B)$ short releases and
$O(P_{\mathrm{elem}}RN_tK_{\mathrm{opt}}B)$ feature differences (see Table~\ref{tab:scaling}).

\begin{table}[t]
\caption{Resources and system boundary of the proposed implementation. The counts
describe the polynomial model; hardware performance depends on the platform. $N$ is
the number of degrees of freedom, $\mathcal E$ the interaction graph of retained
pairs, $P_{\mathrm{elem}}=4N+4|\mathcal E|$ the number of coefficients, $N_t$ the
number of stored noise times, $B$ the batch size, and $R$ the number of repeated CD1
releases. Gradient costs are for one batch at one noise time.}
\label{tab:scaling}
\small
\begin{tabular*}{\textwidth}{@{\extracolsep\fill}
>{\raggedright\arraybackslash}p{0.20\textwidth}
>{\raggedright\arraybackslash}p{0.20\textwidth}
>{\raggedright\arraybackslash}p{0.50\textwidth}}
\toprule
Resource & Nominal scaling & Notes \\
\midrule
Physical state and edges & $O(N+|\mathcal E|)$ & $N$ state elements and a
platform-dependent implementation of each retained interaction. \\
Coefficients & $P_{\mathrm{elem}}=4N+4|\mathcal E|$ & Four on-site coefficients per coordinate and
four coupling coefficients per retained pair. \\
Schedule storage & $O(P_{\mathrm{elem}}N_t)$ & Stored knots for temporal replay;
precision and interpolation add platform-dependent cost. \\
Force-matching gradient & $B$ clamps and $O(BP_{\mathrm{elem}})$ local
contributions & Each clamp provides $N$ force-component values and one contribution
for every coefficient. \\
CD1 gradient & $BR$ short releases and $O(BRP_{\mathrm{elem}})$ local feature
differences & $R$ repetitions reduce variance; the release duration sets the
finite-time bias. \\
Temporal replay & $O(P_{\mathrm{elem}}N_t)$ values & Average delivery rate
$\sim P_{\mathrm{elem}}N_t/\tau$; multiplexing trades channels for bandwidth. \\
Time-to-space mapping & $O(N_t)$ stages and $O(P_{\mathrm{elem}}N_t)$ programmed
elements & Trades control bandwidth for device area and transport. \\
Peripheral cost & Platform dependent & State preparation, sensing, schedule memory,
and controller operation; quantifying these requires a specific platform. \\
\botrule
\end{tabular*}
\end{table}

Electrical circuits provide one possible route. Coupled noisy RLC circuits have
experimentally realized programmable linear Langevin dynamics
\cite{melanson2025thermodynamic}, while feedback-controlled electrical circuits have
realized time-dependent nonlinear effective potentials in one dimension
\cite{chang2021generation}. Networked quartic Langevin models have also been studied
numerically \cite{whitelam2026nonlinear}. A circuit for the present model would need
to combine these capabilities. For example, a node voltage could represent $x_n$,
and a nonlinear current--voltage element could implement the on-site potential. The
four pair features would require platform-specific nonlinear couplers or active
feedback.
Johnson or injected noise supplies the bath. With circuit-specific calibration, the
holding current at a clamped node can estimate one force component. The nonlinear
coefficients, mobility, and noise must all be calibrated.

Diffusion requires the coefficients to change during generation. The cost of
replaying this schedule lies within the system boundary. A spatial device could
instead encode the schedule along a propagation
axis, trading runtime control for device area. It would still need to implement
the stochastic nonlinear dynamics and the trained longitudinal profile. End-to-end
efficiency depends on the specific hardware implementation.

\section{Discussion}\label{sec:discussion}

The time-indexed score field of a diffusion generator can be represented by a driven
conservative energy and trained using quantities accessible to physical measurement.
The reported tests use simulated forces. At positive noise times, the direct
force-matching estimator is unbiased for the score-matching gradient. The
short-time contrastive rule
replaces force sensing with repeated releases, at the cost of finite-time bias and
additional sampling variance. As in standard score-model training, each forward
noise time has a separate objective. Here, each objective can be optimized from
parameter-local quantities without differentiating a numerical device model.
Parameter locality means that each update uses only the coordinate and total-force
components appearing in its energy feature. Preparing a sample still requires
clamping the full state, and CD1 evolves the full coupled system.

A stochastic-circuit implementation still requires an open-loop controller to store
and replay the learned schedule. An alternative strategy is to explore a spatial
encoding of this schedule along a waveguide propagation axis. Longitudinal
modulation and Kerr nonlinearity have been demonstrated in waveguide arrays
\cite{juergensen2021quantized}. A mapping to the dynamics used here remains open.

In the two-dimensional example, the learned finite-time protocol produces samples
in both modes, whereas the persistent equilibrium chain remains in one basin over
the same total simulated time. The nonequilibrium protocol therefore avoids the
slow inter-basin mixing encountered by the static chain for this initialization and
runtime.
The MNIST experiment applies the construction to 784 continuous degrees of freedom
and produces recognizable but blurred digit patterns.

Experimental work should test the rule in a programmable stochastic circuit and
measure controller, sensing, and state-preparation costs. 
A latent-space formulation could decouple the number of oscillators from the data
dimension \cite{vahdat2021score}, but would require a digital encoder and decoder.
Related associative-memory models use a separation of time scales to integrate out
fast hidden variables and obtain effective higher-order interactions
\cite{krotov2020large}. We are currently exploring this idea to extend the parameter-local rule to hidden oscillators. 

\section{Methods}\label{sec:methods}

\subsection{Empirical training distributions}

At each positive noise time, we sample directly from
$p_t^{\mathcal D}$ in Eq.~\eqref{eq:empirical-smoothed}. This distribution is
continuous for $t>0$.
The endpoint is handled implicitly by continuing the sequential optimization to
$t=0$. Here $p_0^{\mathcal D}$ is atomic and has no ordinary density score. The
finite polynomial model cannot represent this singular limit and instead gives a
smooth interpolation between the training samples~\cite{kamb2024analytic}. The exact score-matching
interpretation therefore applies only for $t>0$.
The MNIST images~\cite{lecun1998gradient} are non-binarized grayscale arrays, which we treat as points in
$\mathbb R^{784}$. 
The data set was standardized for training.

\subsection{Force-matching derivation}

We substitute $\vect s_{\vect\theta}=-\nabla\Ehat_{\vect\theta}/\kBT$ into the
Hyv\"arinen objective. Since $\phi_i=\partial_{\theta_i}\Ehat$ and parameter and
state derivatives commute,
\begin{align}
 \partial_{\theta_i}\nabla\!\cdot\vect s_{\vect\theta}
 &=-\Delta\phi_i/\kBT,\\
 \partial_{\theta_i}\tfrac12\lVert\vect s_{\vect\theta}\rVert^2
 &=\vect s_{\vect\theta}\cdot(-\nabla\phi_i/\kBT)
 =-\nabla\phi_i\cdot\Fhat_{\vect\theta}/(\kBT)^2,
\end{align}
which gives Eq.~\eqref{eq:local-gradient}. The derivation does not differentiate
through the reverse-time trajectory because a separate objective is optimized at
each forward noise time.


\subsection{Training protocol}

We initialized $\alpha_n=-1$, $\beta_n=1$, and $\gamma_n=1$ for every site. We
initialized the linear field coefficients $f_n^{\mathrm{ext}}$ and all couplings to
zero.
During optimization, we constrained $\gamma_n$ to exceed $0.001$ to prevent an
unbounded polynomial potential. The coefficients reported in this Article
parameterize the learned score energy $\Ehat$, not the reverse energy $U$.

We optimized sequentially over noise times: the solution at one time initialized
the next. We interpolated the resulting discrete coefficient schedules for
inference. The horizon is $\tau=5$. The mixture experiments use
$N_t=15$ time points, with $\kBT=0.1$ for force matching and $\kBT=0.01$ for
CD1. The MNIST experiment uses $N_t=100$ and $\kBT=0.02$. These values were
selected through a coarse hyperparameter search, and all training and simulations
were run on one NVIDIA L40 GPU.

\authorcheck{Add the CD1 release duration $\delta$ and repeat count $R$.}


\subsection{Use of generative AI in manuscript preparation}

During preparation of this manuscript, the authors used OpenAI Codex for language
editing, structural suggestions, literature searches, and mathematical consistency
checks. The authors reviewed and independently verified all AI-assisted text,
references, and equations and take full responsibility for the content.

\backmatter

\bmhead{Data availability}

The MNIST dataset is publicly available. The numerical data and learned coefficient
schedules underlying the figures will be made available in the accompanying
repository.
\authorcheck{Insert the repository or archive DOI.}


\bmhead{Code availability}

The code, configuration files, and scripts used to generate all numerical results
will be made available in a public repository upon publication and to editors and
referees during review.
\authorcheck{Insert the repository or archive DOI.}

\bmhead{Acknowledgements}

We thank Cindy Zhang for helpful discussions. This research was funded in whole or
in part by the Swiss National Science Foundation (SNSF) through C.B.'s
Postdoc.Mobility fellowship P500PT~217673/1. For the purpose of open access, a CC BY
public copyright licence is applied to any author accepted manuscript (AAM) version
arising from this submission.
This work was performed in part at the Aspen Center for Physics, which is supported
by National Science Foundation grant
PHY-2210452. M.S.-G. was funded by the European Union. Views and opinions expressed
are those of the authors only and do not necessarily reflect those of the European
Union or the European Research Council Executive Agency; neither the European Union
nor the granting authority can be held responsible for them. This work was supported
by ERC grant 101040117 (INFOPASS) and NSF grant OAC-2118201. The funders had no
role in study design, data collection, analysis, interpretation, or preparation of
the manuscript.

\bmhead{Author contributions}

\authorcheck{Insert the author-approved contribution statement using C.B., G.R.,
M.S.-G., and R.P.A.}

\bmhead{Competing interests}

The authors declare no financial or non-financial competing interests.

\bibliography{apssamp}

\end{document}
